\section[APP relevée et triplet spectral en dimension 11]{Relèvement opératoriel d'APP et triplet spectral de l'interface supersymétrique en dimension \texorpdfstring{\(11\)}{11}}
\label{sec:sucesor83-terna-espectral-teoria-m}

Les écrans \(11\to10\), la réduction isotrope \(26\to24\) et l'
involution entre impulsion et enroulement fournissent le support algébrique
exact de l'interface. La dynamique physique ajoute un espace des états, un
hamiltonien et une supercharge. Ces données sont réunies dans le triplet
\[
 \mathfrak M
 =\bigl(\mathcal H_{\mathrm{phys}},H,Q\bigr),
\]
où \(\mathcal H_{\mathrm{phys}}\) est un espace de Hilbert
\(\mathbb Z_2\)-gradué,
\[
 \mathcal H_{\mathrm{phys}}
 =\mathcal H_+\oplus\mathcal H_-,
\]
\(H\) est un opérateur autoadjoint non négatif compatible avec les transports
et \(Q\) est un opérateur impair. La relation de supercharge est
\begin{equation}
 Q^2=H,
 \qquad
 [H,Q]=0.
 \label{eq:sucesor83-supercarga-cuadrado-hamiltoniano}
\end{equation}
L'égalité implique que \(Q\) échange les secteurs gradués et que l'
hamiltonien préserve chacun d'eux.

\subsection{Relèvement opératoriel de l'état APP--TRIT--TPK}

L'APP enrichie conserve simultanément les restes et les quotients des
feuillets additif et multiplicatif :
\[
 \mathcal A_9^\sharp
 =\bigl(R^+,Q^+;R^\times,Q^\times\bigr).
\]
TRIT ajoute l'orientation temporelle et le Topological Phase Kernel conserve
la phase, la route, la retenue, la frontière et la mémoire. Soit
\(\mathcal H_{\mathrm{HMT}}\) l'espace linéaire engendré par ces états
enrichis. Une réalisation opératorielle compatible est une application
\[
 W:\mathcal H_{\mathrm{HMT}}
 \longrightarrow\mathcal H_{\mathrm{phys}}
\]
qui préserve la graduation, transporte les observables de reste et de quotient
et entrelace l'involution bidirectionnelle avec la dualité de rayon. Pour les
opérateurs définis dans les deux domaines, la compatibilité exige
\[
 W\,U_T^{\mathrm{HMT}}
 =U_T^{\mathrm{phys}}\,W,
 \qquad
 W\,H_{\mathrm{HMT}}
 =H\,W.
\]
Cette flèche n'identifie pas par décret les deux espaces : elle spécifie la donnée
que la généalogie doit préserver lors du passage de la construction discrète à sa
réalisation physique.

\subsection{Compactification et codomaine physique undécadimensionnel}

La réalisation en dimension \(11\) a pour codomaine un espace
\[
 \mathcal X_{11}^{\mathrm{phys}}
\]
doté d'une métrique \(g_{MN}\), d'une \(3\)-forme \(C_{MNP}\), d'un gravitino
\(\psi_M\), d'une action et de transformations de supersymétrie. La couture
\(\ell_K\) produit la décomposition exacte
\[
 V_{11}=\ell_K\oplus V_{10},
\]
et la compactification doit induire sur
\(\mathcal H_{\mathrm{phys}}\) la décomposition spectrale correspondante.
La dualité \(T\) exige en outre
\[
 U_T H(R)U_T^{-1}=H(R^{-1}).
\]
La surface d'univers de la corde et les volumes d'univers des branes
apportent alors les réalisations étendues du même antécédent
généalogique.

\subsection{Compatibilité spectrale de \((\mathcal H_{\mathrm{phys}},H,\mathcal Q)\) avec \(\mathcal Q^2=H\), dualité T et compactification \(11\to10\)}

Une réalisation satisfait l'interface lorsqu'elle conserve conjointement :
\begin{enumerate}
\item la graduation de \(\mathcal H_{\mathrm{phys}}\) et l'identité
      \(Q^2=H\) ;
\item l'action de Polyakov et les conditions aux limites de la surface
      d'univers ;
\item la polarisation \(11\to10\) et la réduction isotrope \(26\to24\) ;
\item l'échange de l'impulsion et de l'enroulement sous \(R\leftrightarrow
      R^{-1}\) ;
\item la descente des champs \(g_{MN}\), \(C_{MNP}\) et \(\psi_M\) vers la
      limite compactifiée ;
\item la mémoire de phase, de route, d'orientation, de retenue et de frontière de l'état
      APP--TRIT--TPK.
\end{enumerate}
Les identités des écrans et l'involution de rayon sont exactes dans leurs
domaines algébriques. Le triplet
\((\mathcal H_{\mathrm{phys}},H,Q)\), l'action et les champs
undécadimensionnels fixent, de manière séparée et explicite, les données de l'
interface physique.
